\chapter{Hosting Organisation and Project Framing}
\label{chap:hosting-organisation}

This internship was carried out within Wavestone, a management and technology
consulting group, and more specifically within its Cyber practice in
Luxembourg. Before turning to the regulatory and technical context that
motivates this project, addressed in Chapter~\ref{chap:context}, it is
necessary to situate the organisation that hosted the work and the mission
that was entrusted within it. This chapter presents Wavestone at the group
level, then the Luxembourg office and its Cyber practice, before describing
the PFE mission itself, the working environment and the constraints that
shaped several technical choices, and finally the link between the company's
initial need and the thesis problem developed in the following chapter.

\section{Wavestone: group overview}
\label{sec:wavestone-group}

Wavestone traces its origins to Solucom, a consulting firm founded in 1990,
which merged with the European activities of Kurt Salmon in January 2016 to
form Wavestone as it is known today. In 2023, the group further merged with
Q\_PERIOR, a leading consulting firm in the German-speaking market,
extending its footprint in Germany, Austria and Switzerland. Wavestone is an
independent consulting firm, listed on Euronext Paris (Compartment B), with
its headquarters at the Tour Franklin in La D\'efense, near Paris.

As of the close of its 2025/26 fiscal year (31 March 2026), the group
employed 6,111 consultants across 17 countries, for a consolidated revenue
of 954.3 million euros \cite{wavestone2026results}. Wavestone describes
itself as holding a leading position in its three main continental European
markets, France, Germany and Switzerland, with growing activities in the
United Kingdom and North America. The group organises its consulting offer
into several practices, including for example IT strategy and CTO advisory,
cybersecurity, data and artificial intelligence, and Industry 4.0 and IoT.
Its client base spans large private companies, including numerous CAC40
groups, as well as public institutions, across sectors such as banking,
insurance, energy, transport and the public sector.

\section{Wavestone Luxembourg and the Cyber practice}
\label{sec:wavestone-luxembourg}

Wavestone has had an office in Luxembourg since 2006. As of the writing of
this report, the office employs around fifty consultants, of whom
approximately twenty work within the Cyber practice.

The Cyber practice covers a range of mission types for its clients,
including crisis management exercises, consulting missions on various
topics where consultants act as project leads, and security audits. Its
clients belong mainly to the financial sector, the insurance sector, and
European institutions, in addition to private-sector organisations more
broadly.

\section{The PFE mission within Wavestone}
\label{sec:pfe-mission}

The mission carried out during this internship is titled ``Cyber Crisis
Management Maturity Self-Assessment Platform''. It was carried out within
the Cyber practice introduced in Section~\ref{sec:wavestone-luxembourg}, a
team of around twenty consultants into which the intern was integrated.

The resulting tool is intended to help the practice better assess how
prepared a given organisation is with respect to the ENISA best practices,
and to adapt cyber crisis management exercises accordingly.

Industrial supervision was provided by Riccardo Bergamaschi, Manager within
the Cyber practice.

\section{Working environment and constraints}
\label{sec:working-environment}

The corporate laptop provided by Wavestone did not grant administrator
rights, which made a virtualised development environment mandatory rather
than optional. The platform and its supporting services were therefore run
inside a virtual machine, a constraint that later shaped several of the
networking and deployment choices detailed in Chapter~\ref{chap:platform}.

Running the AI component entirely on local infrastructure, through Ollama,
was a design choice made independently, rather than a requirement dictated
in writing by Wavestone, once the nature of the data was considered: an
assessment is effectively a map of an organisation's cyber weaknesses, and
the NIS2
directive, discussed in Chapter~\ref{chap:context}, sets clear expectations
around data sovereignty. This choice was formalised as one of the five
inviolable principles guiding the project (P2, local-first, see
Chapter~\ref{chap:method}).

Day-to-day work at the Luxembourg office took place in a mix of French and
English, the language used depending on the colleagues involved and on the
project at hand.

\section{From mission to thesis problem}
\label{sec:mission-to-thesis}

The mission described in the previous sections had a well-defined starting
point: the Cyber practice needed a way to give substance to crisis-readiness
conversations with clients, but had no structured tool to do so. ENISA's
2024 study, ``Best Practices for Cyber Crisis Management'', identifies
sixteen best practices spread across the four phases of a crisis,
prevention, preparedness, response and recovery, yet it provides no
mechanism to measure how well an organisation actually applies them. At the
same time, the NIS2 directive puts organisations under growing regulatory
pressure to demonstrate a level of crisis-management maturity that goes
beyond a simple declarative claim. Existing self-assessment tools, such as
the CISA Cyber Resilience Review or the NIST Cybersecurity Framework
self-assessments, were designed independently of ENISA's structure and of
NIS2's specific obligations, and do not map cleanly onto either. Turning the
ENISA framework into an operational, auditable self-assessment platform,
capable of producing a defensible score and AI-assisted recommendations
while keeping client data within the organisation's own infrastructure, is
the concrete engineering problem this mission set out to address.
Chapter~\ref{chap:context} develops the regulatory and technical context
behind this problem in more detail, and Chapter~\ref{chap:method} then
presents the methodological principles used to address it.